\section{\textbf{Evaluating Superintelligence Through Knowledge Invention}}
\label{sec:evaluating_superintelligence}

The Knowledge-Invention Proposition creates an unusual evaluation problem.
Ordinary benchmarks specify tasks, answers, and scoring rules in advance. An
unknown unknown cannot be specified without ceasing to be unknown to the
evaluator. Consequently, superintelligence cannot be established by adding a
more difficult fixed benchmark to an existing test suite. It must be evaluated
through an auditable process that determines whether a generally intelligent
machine originated, validated, and transferred a consequential contribution
beyond the human epistemic frontier.

This section proposes such a process. It is an evaluation framework rather
than a single numerical benchmark. Its purpose is to make claims of machine
knowledge invention testable while avoiding the circular requirement that
human evaluators identify the missing knowledge beforehand. The framework
adopts the testing, evaluation, verification, and validation distinction used
in established AI-assurance practice \cite{tabassi2023airmf}, but applies it to
the epistemic and attributional conditions developed in this paper.

\subsection{\textbf{Why Fixed-Task Evaluation Is Insufficient}}

A fixed task can measure whether a system solves a problem humans already
recognize. It may reveal reasoning ability, transfer, planning, or performance
beyond individual experts. It cannot by itself establish that the system
identified knowledge absent from the recognized human research agenda. Once
the evaluator supplies the question, target, candidate space, and success
criterion, the central unknown-unknown operation has already been performed by
humans.

Public benchmarks create an additional problem: exposure of evaluation items
or their close derivatives can inflate measured capability. Data contamination
has therefore become a material threat to valid language-model evaluation
\cite{sainz2023contamination}. For knowledge invention, contamination includes
not only exact answer leakage but also prior exposure to the decisive
hypothesis, relation, experimental design, or unpublished human result.

Let $B$ be a fixed benchmark and $s_B(x)$ the score of system $x$. Then

\begin{equation}
\begin{aligned}
s_B(x) > s_B(H)
\not\!\Longrightarrow
\operatorname{KI}_{x}=1,
\end{aligned}
\label{eq:benchmark_not_invention}
\end{equation}

where $H$ denotes the relevant human comparison group. Equation~
\eqref{eq:benchmark_not_invention} does not diminish benchmark value; it limits
the inference that may be drawn from benchmark success. A benchmark may
establish component competence, but knowledge invention requires evidence
about novelty, origination, consequence, and validation.

\subsection{\textbf{Evaluation Object and Frozen Record}}

The unit of evaluation is not an isolated model response. It is an invention
episode

\begin{equation}
\begin{aligned}
\mathcal{E}
=(&x,\theta,D,T,H,A,\tau,\kappa),
\end{aligned}
\label{eq:invention_episode}
\end{equation}

where $x$ is the evaluated system, $\theta$ its fixed model and policy
configuration, $D$ the accessible data, $T$ the available tools, $H$ all human
interventions, $A$ the machine action trace, $\tau$ the evaluation interval,
and $\kappa$ the resulting candidate contribution. The record must be frozen
before the candidate is assessed so that later success cannot silently alter
the apparent degree of machine autonomy.

The frozen record should contain model versions, system and user instructions,
retrieval results, code, tool calls, intermediate hypotheses, rejected paths,
experimental decisions, human messages, timestamps, and generated artifacts.
This requirement follows the broader scientific need for transparent and
reproducible reporting \cite{munafo2017reproducible}, but it also serves a
distinct purpose: reconstructing who or what supplied the decisive epistemic
content.

\subsection{\textbf{Stage I: Frontier and Agenda Verification}}

The first stage evaluates whether the candidate lies beyond both known
knowledge and the recognized research agenda. Because no search can prove
absolute absence, novelty is assessed through convergent evidence. Relevant
evidence includes time-stamped scholarly and patent searches, structured
database searches, expert review, comparison with negative or unpublished
results when available, and semantic searches for equivalent formulations.

Two boundaries must be evaluated separately:

\begin{equation}
\begin{aligned}
F_K(\kappa)&=1
&&\text{iff no materially equivalent contribution is found},\\
F_Q(q_\kappa)&=1
&&\text{iff no materially equivalent research question is found},
\end{aligned}
\label{eq:frontier_tests}
\end{equation}

where $F_K$ is the knowledge-frontier test and $F_Q$ the agenda-frontier test.
A candidate may pass $F_K$ while failing $F_Q$: it may solve an open problem
already identified by humans. Such a result can be an important machine
discovery, but it does not satisfy the stronger unknown-unknown condition.

The novelty panel should include independent domain specialists and at least
one reviewer assigned to search actively for defeating evidence. Review must
be blinded to system identity where feasible. The panel should report the
search scope and a calibrated conclusion rather than claiming logical proof of
absence.

\subsection{\textbf{Stage II: Unknown-Unknown Origination}}

The second stage asks how the research object entered the workflow. The key
question is not whether a human touched the process, but whether human input
contained the decisive question, hypothesis, representation, or evaluation
criterion. A machine may use human-built instruments and still originate the
invention; conversely, an autonomous search can remain fully enclosed within a
human-specified epistemic frame.

Let $I_H$ denote the semantic information supplied by humans before the
candidate appeared, and let $d$ denote the decisive epistemic content of
$\kappa$. Define

\begin{equation}
\begin{aligned}
U_x(\kappa)=1
\Longleftrightarrow{}&
d\not\subseteq I_H\\
&\land
\operatorname{TraceOrigin}(d,A)=x\\
&\land
\operatorname{Counterfactual}(W\setminus d)
\neq \kappa,
\end{aligned}
\label{eq:unknown_unknown_origination}
\end{equation}

where $W$ is the complete workflow. The first condition excludes substantive
human specification, the second requires trace evidence of machine
origination, and the third requires the machine contribution to be causally
decisive. Ablations, replay experiments, independent trace annotation, and
comparison against matched human-guided runs can strengthen this assessment.

\subsection{\textbf{Stage III: Independent Validation}}

Novelty and origination do not establish truth. The candidate must next face a
validation procedure chosen for its epistemic type: formal proof checking for
a theorem or algorithm, controlled experiment for a causal or empirical claim,
prospective prediction for a scientific model, synthesis and measurement for
a proposed material, or deployment under prespecified conditions for an
engineering method. Autonomous laboratories already demonstrate how planning,
experimentation, analysis, and revision can be connected in a closed research
loop \cite{king2004robot,king2009automation,boiko2023autonomous}; the present
framework adds independence between origination and acceptance.

The originating system must not be the sole validator. Let
$\mathcal{V}=\{v_1,\ldots,v_n\}$ be independent validation channels. Then

\begin{equation}
\begin{aligned}
V(\kappa)=1
\Longleftrightarrow{}&
\sum_{j=1}^{n}w_jv_j(\kappa)\geq\gamma\\
&\land
\operatorname{Independence}(\mathcal{V},x)\geq\delta\\
&\land
\operatorname{Precommit}(\mathcal{V})=1,
\end{aligned}
\label{eq:independent_validation}
\end{equation}

where $w_j$ reflects evidentiary strength, $\gamma$ is the domain-specific
acceptance threshold, and $\delta$ is a minimum independence requirement.
Precommitment means that the principal acceptance criteria are recorded before
validation results are known. Failed or mixed validation must remain in the
record rather than being removed retrospectively.

\subsection{\textbf{Stage IV: Consequentiality}}

A valid novelty may still be too narrow or trivial to demonstrate the form of
superintelligence proposed here. Consequentiality concerns whether the
contribution changes what can be explained, predicted, proved, designed, or
investigated. It must be assessed against explicit counterfactual baselines,
not excitement surrounding an AI-generated result.

We represent consequentiality as a multidimensional judgment:

\begin{equation}
\begin{aligned}
C(\kappa)=\Omega(&\Delta E,\Delta P,\Delta R,
\Delta G,\Delta T),
\end{aligned}
\label{eq:consequentiality_vector}
\end{equation}

where $\Delta E$ is explanatory gain, $\Delta P$ predictive gain, $\Delta R$
research enablement, $\Delta G$ generality across contexts, and $\Delta T$
transformative practical effect. Not every contribution must score highly on
all dimensions. It must, however, produce a material advance on at least one
prespecified dimension without losing validity on the others relevant to its
claim. Atypical combinations can generate unusually influential science
\cite{uzzi2013atypical}, but unusualness alone is not consequentiality.

\subsection{\textbf{Stage V: Replication and Capability Testing}}

An accepted invention establishes an event; it does not necessarily establish
a repeatable capability. The evaluated system should therefore be tested over
multiple invention episodes with different data environments, tools, domains,
and human teams. Negative attempts must be counted so that selection of one
successful episode cannot conceal a low or chance-level success rate.

For $m$ registered episodes, define

\begin{equation}
\begin{aligned}
\widehat{p}_{\mathrm{KI}}(x)
&=\frac{1}{m}\sum_{i=1}^{m}
\mathbf{1}[F_K\land F_Q\land U_x\land V\land C]_i,\\
R_x&=\operatorname{ReplicateRate}(x),\\
D_x&=\operatorname{DomainDiversity}(x).
\end{aligned}
\label{eq:capability_profile}
\end{equation}

The resulting capability profile is

\begin{equation}
\begin{aligned}
\Pi_x=
(\widehat{p}_{\mathrm{KI}},R_x,D_x,
\mathcal{A}_x,\mathcal{U}_x),
\end{aligned}
\label{eq:evaluation_profile}
\end{equation}

where $\mathcal{A}_x$ aggregates attribution strength and $\mathcal{U}_x$
records uncertainty. Reporting the profile is preferable to collapsing every
dimension into one leaderboard score. It preserves the distinction between a
system that repeatedly invents within one scientific field and a system that
transfers the capability across domains.

\subsection{\textbf{Decision Rule and Falsification}}

For each candidate, the evaluation panel returns one of four outcomes:
rejected, unresolved, validated machine discovery, or validated machine
knowledge invention. A contribution qualifies for the final category only if
all gates are passed:

\begin{equation}
\begin{aligned}
\operatorname{KI}_{x}(\kappa)=1
\Longleftrightarrow{}&
F_K(\kappa)=1\\
&\land F_Q(q_\kappa)=1\\
&\land U_x(\kappa)=1\\
&\land V(\kappa)=1\\
&\land C(\kappa)=1.
\end{aligned}
\label{eq:evaluation_gates}
\end{equation}

The claim is falsified for a candidate if a materially equivalent prior result
or research question is found, the decisive content is traced to human input
or inaccessible leakage, independent validation fails, or the claimed
consequence disappears under the stated baseline. A system-level claim is
weakened or defeated when success cannot be replicated, depends on one narrow
environment, or vanishes under prospective evaluation.

This protocol deliberately makes superintelligence difficult to claim. High
performance on known tasks, autonomous execution, broad language competence,
and even a major solution to a recognized open problem remain important forms
of machine intelligence. The stronger designation is reserved for a generally
intelligent system that repeatedly originates and validates consequential
knowledge beyond both the known human frontier and the questions humanity had
already learned to ask.
